% ============================================================
% ch19.tex —— 第 19 章 e：自己会求导的函数
% ============================================================
\chapter{\texorpdfstring{$\mathrm{e}$}{e}：自己会求导的函数}

微积分工具箱里还缺一位主角。它藏在一条银行条款里，管着全球的复利、人口增长和放射性衰变，还将作为篇四素数普查的标尺。它的数值是 $2.718281828\cdots$，名字只有一个字母：$\mathrm{e}$。这一章我们把它发明出来——并且证明它配得上"自然"二字。

\section{银行家的世纪难题}

\begin{kaiwei}
存 $1$ 元，年利率 $100\%$（夸张的假设，方便看清结构）。一年后连本带息是 $2$ 元。现在银行提供"复利升级"：半年结一次息（年利率不变，每次结一半），每季度结一次、每月结一次、每天结一次……甚至\textbf{连续计息}。结算次数无限多时，一年后能拿多少？能靠多结息\textbf{无限暴富}吗？
\end{kaiwei}

冷静地算。每次结息，本金乘上 $\left(1 + \tfrac1n\right)$（一年结 $n$ 次，每次利率 $\tfrac1n$），一年下来连乘 $n$ 次：

\begin{center}
\begin{tabular}{ccc}
\toprule
结算频率 $n$ & 计算 $\left(1+\dfrac1n\right)^n$ & 一年后的倍数 \\
\midrule
$1$（年结） & $(1+1)^1$ & $2$ \\
$2$（半年结） & $(1+\tfrac12)^2 = 1.5^2$ & $2.25$ \\
$4$（季结） & $(1.25)^4$ & $2.44140625$ \\
$12$（月结） & $\left(1+\tfrac1{12}\right)^{12}$ & $2.613035\cdots$ \\
$365$（日结） & $\left(1+\tfrac1{365}\right)^{365}$ & $2.714567\cdots$ \\
$8760$（时结） & $\left(1+\tfrac1{8760}\right)^{8760}$ & $2.718127\cdots$ \\
$525600$（分结） & $\left(1+\tfrac{1}{525600}\right)^{525600}$ & $2.7182792\cdots$ \\
\bottomrule
\end{tabular}
\end{center}

涨得越来越慢：从年到半年，白赚 $0.25$；从月到日（频率增三十倍），只多赚 $0.10$；从时到分（频率增六十倍），只多赚 $0.0001$。数字明显在逼近某个\textbf{封顶}——又是一场落脚点（第 14 章的比赛规则第 $N$ 次上场）。欧拉 1748 年在《无穷分析引论》里给它发了正式户口：
\[
\mathrm{e} := \lim_{n\to\infty}\left(1+\frac1n\right)^n = 2.718281828459\cdots
\]
\textbf{连续复利不会让你无限暴富}：$1$ 元最多变成 $\mathrm{e} \approx 2.718$ 元。结算越频繁、回报越接近上限，但\textbf{上限纹丝不动}——极限思想第一次出手，管住了金融学的贪婪。（你存的"日计息"理财，实付 $2.7146$ 元档；想拿到 $2.71828$，得"每秒结息"再算上四百万年——差距小到银行懒得区分，这就是"连续复利"作为工程近似的底气。）

\section{给 $\mathrm{e}^x$ 拍 X 光片：唯一求导不变的函数}

$\mathrm{e}$ 的真正威力在函数 $\mathrm{e}^x$ 身上。用第 15 章的机器给它做体检：
\[
\frac{\mathrm{e}^{x+h} - \mathrm{e}^x}{h} = \frac{\mathrm{e}^x\cdot \mathrm{e}^h - \mathrm{e}^x}{h} = \mathrm{e}^x\cdot\frac{\mathrm{e}^h - 1}{h} .
\]
（$\mathrm{e}^{x+h} = \mathrm{e}^x\mathrm{e}^h$，指数相加 $=$ 底数相乘，初中的同底数幂法则对无理指数照常成立。）前一个因子 $\mathrm{e}^x$ 与 $h$ 无关，原样带走；成败押在后一个因子 $\dfrac{\mathrm{e}^h - 1}{h}$ 上——它只含 $h$。查它的数值：

\begin{center}
\begin{tabular}{cccccc}
\toprule
$h$ & $0.1$ & $0.01$ & $0.001$ & $0.0001$ & $\to 0$ \\
\midrule
$\dfrac{\mathrm{e}^h - 1}{h}$ & $1.0517$ & $1.0050$ & $1.0005$ & $1.00005$ & $\to 1$ \\
\bottomrule
\end{tabular}
\end{center}

落脚点是 $1$（这不是巧合，而是 $\mathrm{e}$ 的定义式的对偶——诚实声明里给出联系）。于是差商的落脚点：
\[
\boxed{\;\left(\mathrm{e}^x\right)' = \mathrm{e}^x\;}
\]
\textbf{求导不改变它}。全宇宙的函数里，"导数等于自己"只此一家的倍数：$(c\,\mathrm{e}^x)' = c\,\mathrm{e}^x$，而 $2\mathrm{e}^x$、$-5\mathrm{e}^x$ 全是 $\mathrm{e}^x$ 的倍数。\textbf{选 $\mathrm{e}$ 当底数，恰好像给指数函数装了免检牌照}：增长率永远等于当前值。

这为什么是"自然的"？看反面：底数 $2$ 的函数 $2^x$，差商提出 $2^x$ 后剩下 $\dfrac{2^h - 1}{h}\to 0.6931\cdots$（不是 $1$），于是 $(2^x)' = 0.693\,2^x$——拖着一个尾巴系数；底数 $10$，尾巴 $2.3026\cdots$；底数 $\mathrm{e}$，尾巴恰好 $1$。\textbf{$\mathrm{e}$ 不是被随机选中的吉祥物，是让"增长率 $=$ 当前值"这句话成立的唯一底数}。自然界"长出来的部分接着长"的过程——人口（资源无限时）、细菌繁殖、放射性衰变（负增长版）、连续复利——全都是这条微分规律的演出，所以它们清一色是 $\mathrm{e}$ 的指数。$\mathrm{e}$ 之"自然"，如同圆之"$\pi$"：不是审美，是结构。

\section{顺手发明对数}

方程 $\mathrm{e}^x = 2$，$x$ 是多少？$x$ 一定存在（$\mathrm{e}^0 = 1 < 2 < \mathrm{e} = \mathrm{e}^1$，而 $\mathrm{e}^x$ 连续上涨，中途必经 $2$——"介值"直觉），但它是一个"没有名字的数"。按序章的发明许可证办：\textbf{给它起名}。
\[
\ln a \;:=\; \text{方程 } \mathrm{e}^x = a \text{ 的解} \qquad (\text{"}a\text{ 的自然对数"}) .
\]
对数不是背出来的，是\textbf{被解方程逼出来的}——和分数（第 10 章）、负数（序章）的诞生同一剧本。$\ln 2 \approx 0.6931$（上面底数 $2$ 的尾巴系数，一会儿解释它为何同一副面孔）、$\ln \mathrm{e} = 1$、$\ln 1 = 0$。

$\ln$ 的运算法则不靠背，\textbf{从指数法则翻译}。指数说 $\mathrm{e}^{u+v} = \mathrm{e}^u\cdot\mathrm{e}^v$；两边取 $\ln$（"取 $\ln$"就是问"这是 $\mathrm{e}$ 的几次方"）：
\[
u + v = \ln\left(\mathrm{e}^u\mathrm{e}^v\right) \quad\Longrightarrow\quad \boxed{\ln(ab) = \ln a + \ln b}
\]
（代入 $a = \mathrm{e}^u, b = \mathrm{e}^v$。）同理 $\ln\dfrac ab = \ln a - \ln b$，$\ln a^n = n\ln a$。

\textbf{乘法降维成加法}。这条式子改写过文明史：纳皮尔 1614 年发明对数（比 $\mathrm{e}$ 被认识还早——他造的是以 $\tfrac{1}{\mathrm{e}}$ 为底的表，冥冥中又撞见 $\mathrm{e}$），开普勒拿对数算出了火星轨道；此后三百年，全世界的科学家、工程师、会计人手一册对数表，把 $437\times892$ 化成"查 $\ln$、相加、查反表"——一次乘法省下的笔算约二十分钟，一生的科研时间多出数年。\textbf{计算尺}（两条对数刻度的尺子，滑动相加实现相乘）是对数时代的标志性工具：阿波罗登月的舱内还备着它，作为计算机的机械备份。

\begin{gongju}
\textbf{工具箱（随取随用，篇四全程依赖）}：
\[
\ln(ab) = \ln a + \ln b,\qquad \ln\frac ab = \ln a - \ln b,\qquad \ln(a^n) = n\ln a,
\]
\[
\ln 1 = 0,\quad \ln \mathrm{e} = 1,\quad \ln 2 \approx 0.6931,\quad \ln 10 \approx 2.3026 .
\]
推论速算：$\ln 8 = 3\ln 2 \approx 2.079$；$\ln 100 = 2\ln10 \approx 4.605$；$\ln 1000 \approx 6.908$；$\ln 0.5 = -\ln 2 \approx -0.693$。
\end{gongju}

\section{级数彩蛋：$\mathrm{e}$ 的无尽加法}

$\mathrm{e}$ 还有第三副面孔——一条收敛飞快的级数（欧拉发现，此处直接展示）：
\[
\mathrm{e} = 1 + 1 + \frac{1}{2!} + \frac{1}{3!} + \frac{1}{4!} + \frac{1}{5!} + \cdots = \sum_{k=0}^{\infty}\frac{1}{k!},
\qquad k! = 1\times2\times\cdots\times k .
\]
手算前八项（$k = 0$ 到 $7$）：
\[
1,\ 2,\ 2.5,\ 2.6667,\ 2.7083,\ 2.7167,\ 2.71806,\ 2.718254,\ \dots
\]
八项咬住小数点后四位。收敛为什么这么快？\textbf{分母是阶乘}——比几何级数的"翻倍"还凶猛地膨胀（$2, 6, 24, 120, 720$……每项至少缩小到前项的几分之一且越缩越狠）。对照第 13 章的预告：\textbf{级数收敛与否、收敛快慢，全看分母的性格}——$+1$ 爬坡的调和级数发散（第 21 章），平方增长的收敛到 $\tfrac{\pi^2}6$，阶乘增长的收敛如闪电。

\paragraph{选读小节：最美丽的公式。}允许虚数单位 $\mathrm{i}$（$\mathrm{i}^2 = -1$，把数轴升级成平面——横轴实数、纵轴虚数）后，欧拉写下：
\[
\mathrm{e}^{\mathrm{i}\theta} = \cos\theta + \mathrm{i}\sin\theta ,
\]
即：\textbf{$\mathrm{e}^{\mathrm{i}t}$ 以角速度 $1$ 沿单位圆逆时针转圈}。为什么？用"变化率"论证（本书的招牌）：设 $f(t) = \mathrm{e}^{\mathrm{i}t}$，链式法则给出 $f'(t) = \mathrm{i}f(t)$——"变化率永远等于当前位置逆时针转 $90^\circ$"（乘 $\mathrm{i}$ 的几何作用正是逆时针转 $90^\circ$）。而"位置永远垂直于半径、速度恒为 $1$"的运动，恰是\textbf{匀速圆周运动}：$f$ 从 $1$ 出发绕单位圈，$t$ 时刻走到角度 $t$ 处，即坐标 $(\cos t, \sin t)$。两式合一，欧拉公式成。令 $\theta = \pi$（转半圈到 $-1$）：
\[
\mathrm{e}^{\mathrm{i}\pi} = -1 \quad\Longleftrightarrow\quad \mathrm{e}^{\mathrm{i}\pi} + 1 = 0 ,
\]
全数学最著名的五个常数 $\mathrm{e}, \mathrm{i}, \pi, 1, 0$ 在一行里会师，数学界称它"最美公式"。\textbf{$\mathrm{e}$ 不只管增长，还管旋转}——增长与旋转，一个数的两种人生。

\begin{dongshou}
\begin{itemize}
  \yisi 用复利表格的算法手算 $(1+\tfrac1{100})^{100}$ 的近似（二项式展开前三项：$1 + 1 + \tfrac{99}{200} \approx 2.485$，对照精确值 $2.7048$——展开要取更多项才准，此题体会"截断误差"）。
  \yier 求 $f(t) = \mathrm{e}^{3t}$ 的导数（链式：$3\mathrm{e}^{3t}$），推广并验证 $(\mathrm{e}^{kt})' = k\,\mathrm{e}^{kt}$。用数量级直觉检查：$k = 3$ 的增长是 $k=1$ 的三倍速，导数也三倍，自洽。
  \yier 用 $\ln$ 法则化简：$\ln 8 + \ln 9 - \ln 6$（$=\ln 12 \approx 2.485$）；$\ln\tfrac{25}{4}$（$=2\ln5 - \ln4$）；再解方程 $\mathrm{e}^{3x} = 5$（两边取 $\ln$：$x = \tfrac{\ln 5}{3} \approx 0.5365$）。
  \yisi 不用计算器判断大小：$\ln 3$ 与 $1$？$\ln 10$ 与 $2.5$？（$\ln 3 < \ln \mathrm{e} = 1$；$\ln 10 \approx 2.3026 < 2.5$。对数题的第一反应：和 $\ln\mathrm{e}=1$、$\ln 10\approx2.3$ 两个锚点比。）
\end{itemize}
\end{dongshou}

\begin{shaonao}
\textbf{三线合流的预演。}

(1) 连续复利翻倍要多久？$\mathrm{e}^t = 2$，$t = \ln2 \approx 0.693$ 年。速算口诀"七二法则"：年利率 $r\%$ 下翻倍年数 $\approx \tfrac{72}{r}$——用 $\mathrm{e}^{rt} = 2$ 验证 $r = 6$：$t = \tfrac{\ln2}{0.06} \approx 11.55$ 年，$72/6 = 12$，吻合到惊人。请推导口诀的来源：$t = \tfrac{\ln 2}{\ln(1+r)} \approx \tfrac{0.693}{r}$（利率小时 $\ln(1+r)\approx r$），$69.3\approx72$（顺手调大一点，补偿整数误差）。

(2) 放射性碳测年：碳-14 半衰期 $5730$ 年，即每过 $5730$ 年剩余 $\tfrac12$。剩余比例模型 $N(t) = \mathrm{e}^{-\lambda t}$，$\lambda = \tfrac{\ln 2}{5730}$。考古测得某遗物碳-14 剩 $30\%$：$\mathrm{e}^{-\lambda t} = 0.3$，$t = \tfrac{\ln(10/3)}{\lambda} \approx \tfrac{1.204}{1.21\times10^{-4}} \approx 9950$ 年。亲手复算这两步。

(3) 对数级数：$\ln 2 = 1 - \tfrac12 + \tfrac13 - \tfrac14 + \cdots$（交错调和级数）。算前八项（$0.6345$）、前十六项（$0.6629$），感受它的\textbf{慢}——与 $\mathrm{e}$ 的级数（八项四位）对照：分母 $+1$ 爬坡的收敛有多煎熬。慢，但确凿奔 $\ln 2 \approx 0.6931$——第 21 章它还要回来当主演。
\end{shaonao}

\begin{chengshi}
三处地毯：其一，"$\left(1+\tfrac1n\right)^n$ 的极限存在"（单调有界原理，大学分析第一课的招牌例题）；其二，"$\tfrac{\mathrm{e}^h-1}{h}\to1$"与定义式的对偶——令 $h = \tfrac1n$，$\left(1+\tfrac1n\right)^n = \left(\left(1+h\right)^{1/h}\right)^h$，括号内奔 $\mathrm{e}$，故 $1+h = \mathrm{e}^h(1+o(1))$，即 $\mathrm{e}^h - 1 \approx h$（严格版需交换两个极限步骤，柯西的功劳）；其三，欧拉公式的"求导论证"默认了复指数的可导性（复分析第一课）。$\mathrm{e}$ 是无理数（欧拉证）且是超越数（埃尔米特 1873 年证，第 12 章的户口本上有它）。
\end{chengshi}

篇三收官。现在，把篇一的老朋友（素数）和篇三的新发动机（$\ln$、$\mathrm{e}$、积分）装进同一台机器——篇四：解析数论，正式开机。你会看到两百年来最不可思议的一次"跨界联姻"。
